\section{梯度流}
规范场与费米子场的梯度流（GF）~\cite{Luscher:2010iy,Luscher:2011bx,Luscher:2013cpa}
为调节 twist-2 算符的短距离奇异性提供了另一条途径。
在流时间 $t=0$ 处，与物理重正化矩阵元的联系通过短流时间展开（SFTX）获得：
先在固定物理流时间 $t$ 处完成强子矩阵元的连续极限~\cite{Luscher:2013vga}。

用带流费米子场 $\chi^r(x,t)$、$\chibar^r(x,t)$ 和规范场 $B_\mu(x,t)$ 写出的
带流 twist-2 算符为
\be
O_n^{rs}(x,t) = \chibar^r(x,t) \gmuopen1 \lrDmu2 \cdots \lrDmuclose{n} \chi^s(x,t)\,,
\label{eq:flowed_t2}
\ee
其中 $\left\{\cdots\right\}$ 表示对所含指标的对称化。
这些算符特别有利之处在于：即使在格点上它们也是乘法重正化的，
且重正化因子只依赖算符的费米子内容，
\be
O_n^{rs}(t) = Z_n O_{n,B}^{rs}(t)\,, \quad Z_n = Z_\chi\,,
\ee
其中 $Z_\chi^{1/2}$ 是带流费米子场的重正化常数~\cite{Luscher:2013cpa}。
对含带流 twist-2 算符的关联函数，流时间 $t$ 为短距离奇异性提供了调节。
除理论裸参数、带流费米子场以及 $t=0$ 处局域场的重正化之外，
关联函数不需要任何额外重正化。
对格点 QCD 应用而言，方便的做法是在一个与调节无关的方案中定义 $Z_\chi$，
这样匹配与重正化便可在同一方案中进行。
具体方案的选择对一般讨论并不重要，
但在计算匹配系数时变得重要。在下述匹配系数的计算中，
我们使用所谓的环场（ringed fields）~\cite{Makino:2014taa}
重正化 twist-2 算符；环场由 SU($N_c$) 规范不变且与调节无关的条件定义：
\be
\left \langle \rchibar_r(x,t) {\overset\leftrightarrow{\slashed{D}}} \rchi_r(x,t) \right\rangle
= - \frac{N_c}{(4 \pi)^2 t^2}\,.
\ee
在维数正规化（$D=4 - 2 \epsilon$）中施加该条件，
导致环场与 $\MS$ 方案重正化场之间的有限重正化
\bea
\chi_r(x,t) &=& \left(8 \pi t\right)^{\epsilon/2} \zeta_\chi^{1/2} \rchi_r(x,t) \nonumber \\
\chibar_r(x,t) &=& \left(8 \pi t\right)^{\epsilon/2} \zeta_\chi^{1/2} \rchibar_r(x,t)\,.
\eea
该有限重正化的单圈结果为~\cite{Makino:2014taa}
\be
\zeta_\chi = 1 - \frac{\gbar^2(\mu)}{(4 \pi)^2}C_F
\left( 3 \log (8 \pi \mu^2 t) - \log (432)\right)\,,
\ee
其中 $\gbar$ 是在重正化标度 $\mu$ 处 $\MS$ 方案重正化的强耦合常数，
$C_F = N_c^2-1/2N_c$。该结果已被推广到 O($\gbar^4$)
阶~\cite{Harlander:2018zpi,Artz:2019bpr}。
$\MS$ 方案中得到的结果可通过把重正化标度 $\mu$ 换成
$\mubar$ 而转换为 $\MSbar$ 方案，两者关系为
$\log \mu^2 = \log \mubar^2 + \gamma_E - \log 4 \pi$。

\subsection{O($a$) 改进}
twist-2 算符强子矩阵元的格点 QCD 计算受截断效应影响。

对非微扰改进的 clover 费米子，
标准 twist-2 算符（$t=0$）的强子矩阵元仍带有 O($a$) 截断效应；
采用 O($a$) 改进的 twist-2 算符定义即可将其消除。
除 $n=2$ 情形外——其相应改进系数已在微扰论单圈水平确定~\cite{Capitani:2000xi}——
连续极限以 O($a$) 离散化效应到达。

若考虑一般 $n$ 的带流 twist-2 算符，
其强子矩阵元只受 O($a m$) 截断效应影响（$m$ 为夸克质量），
不需要任何 twist-2 算符专属的附加改进系数。
在下文方法应用的讨论中，我们还考虑带流矩阵元的比值，
此时 $Z_\chi$ 与 O($am$) 截断效应同时被消去。
总而言之，无论采用何种格点方案，
使用带流 twist-2 算符不仅显著简化了重正化模式，
还改善了向连续极限过渡时相对格距的参数化标度行为。
